\documentstyle[fleqn,figfonts,amstex,epsf]{article}
\setlength{\textwidth } {16cm} 
\setlength{\oddsidemargin} {1cm} 
\setlength{\evensidemargin} {1cm} 
\setlength{\topmargin} {1cm} 
\setlength{\headheight}{0cm}
\setlength{\headsep}{0cm}
\setlength{\textheight} {22.5cm}
\setlength{\parindent}{10mm}
\title{Some properties of the symplectic Lie algebra}
\author{R.Coleman\\Laboratoire LMC-IMAG,\\Tour IRMA, 51 rue des Math\'ematiques,
\\Domaine Universitaire BP 53\\38041 Grenoble Cedex 9, France.}
\newcommand{\idV}{\operatorname{id_V}}
\newcommand{\icp}{\text i}
\newcommand{\gen}{\operatorname{gen}}
\newcommand{\eig}{\operatorname{eig}}
\begin{document}
\maketitle

In this talk we give a brief introduction to the symplectic Lie algebra, 
which is composed of $2n\times 2n$ matrices $A$ satisfying the property
$$
A^tJ + JA = 0 
$$
where
$$
J = J_{2n} = \left\lgroup \begin{array}{cc}
		0 & I_n\\
		-I_n & 0
		\end{array} \right\rgroup
$$
In particular, we will give an explicit proof of the simplicity of this Lie algebra, 
something which is often alluded to, but not proved.\\
From this property we will draw some conclusions about the symplectic group, the 
Lie group of $2n\times 2n$ matrices $T$ with the property
$$
T^tJT = J
$$
Next we will look at the Lie algebra of smooth functions defined 
on an open subset $O$ of ${\Bbb R}^{2n}$ and show how the simplicity 
of the symplectic Lie algebra raises questions concerning this Lie algebra.\\

\bibliographystyle{plain}
\bibliography{biblio}
\end{document}



 
 







 





--PART-BOUNDARY=.1990505155610.ZM18327.imag.fr--


